% Owner: Member C (Eesha Irfan); mapping assumed from proposal order.
\chapter{Implementation and Test Cases}
For each chapter, provide a paragraph of introduction and at the end, a paragraph of conclusions. (Not the programming code but the algorithmic and procedural details especially related to the hidden/ backend algorithms that are not covered in the design)
\section{Implementation}
Whatever implementation that you have done so far, please elaborate here.
Give clear details of the algorithms that were implemented along with the platform and the APIs which were used. 
\textcolor{red}{NOTE: For FYP-1, this chapter can be changed to description of prototype developed.}
\subsection{Implementation of First Component/Algorithm}
Write the implementation of the first component of your system here.
\section{Test case Design and description}
\textcolor{red}{This section will be added in FYP-II.} Summarize the common attributes of test cases. This may include input constraints that must be true for every input in the set of associated test cases, any shared environmental needs, any shared special procedural requirements, and any shared case dependencies. The following scheme is recommended for describing test cases in detail. The Sample test case is provided in table \ref{tab:3}.
\begin{table}[!ht]
\caption{Sample Test case No.1}

\centering
\begin{tabular}{|llll|}
\hline
\multicolumn{4}{|c|}{\textbf{\textless{}Software Component Name\textgreater{}}} \\ \hline
\multicolumn{4}{|c|}{\textbf{\textless{}Reference\textgreater{}}} \\ \hline
\multicolumn{1}{|l|}{\textbf{Test Case ID:}} &
  \multicolumn{1}{l|}{\textit{Reference Number}} &
  \multicolumn{1}{l|}{\textbf{QA Test Engineer:}} &
  \textit{Name of personnel} \\ \hline
\multicolumn{1}{|l|}{\textbf{Test case Version:}} &
  \multicolumn{1}{l|}{\textit{Version number}} &
  \multicolumn{1}{l|}{\textbf{Reviewed By:}} &
  \textit{Testing Team lead} \\ \hline
\multicolumn{1}{|l|}{\textbf{Test Date:}} &
  \multicolumn{1}{l|}{\textit{Date}} &
  \multicolumn{1}{l|}{\textbf{\begin{tabular}[c]{@{}l@{}}Use Case \\ Reference(s):\end{tabular}}} &
  \textit{Relation to use cases} \\ \hline
\multicolumn{1}{|l|}{\textbf{Revision History:}} &
  \multicolumn{3}{l|}{\textit{Refer to previous test   case identity (if any)}} \\ \hline
\multicolumn{1}{|l|}{\textbf{Objective:}} &
  \multicolumn{3}{l|}{\textit{Need and   scope of the testing}} \\ \hline
\multicolumn{1}{|l|}{\textbf{\begin{tabular}[c]{@{}l@{}}Product/Ver/\\ Module:\end{tabular}}} &
  \multicolumn{3}{l|}{\textit{Refer to overall system being built and the place of this test case in it.}} \\ \hline
\multicolumn{1}{|l|}{\textbf{Environment:}} &
  \multicolumn{3}{l|}{\textit{Necessary and desired properties of the test environment.}} \\ \hline
\multicolumn{1}{|l|}{\textbf{Assumptions:}} &
  \multicolumn{3}{l|}{\textit{Assumptions that might affect the testing process.}} \\ \hline
\multicolumn{1}{|l|}{\textbf{Pre-Requisite:}} &
  \multicolumn{3}{l|}{\textit{Necessary condition that  needs to be fulfilled before the test case.}} \\ \hline
\multicolumn{1}{|l|}{\textbf{Step No.}} &
  \multicolumn{1}{l|}{\textbf{Execution   description}} &
  \multicolumn{2}{l|}{\textbf{Procedure result}} \\ \hline
\multicolumn{1}{|l|}{} &
  \multicolumn{1}{l|}{\textit{Events being tested.}} &
  \multicolumn{2}{l|}{\textit{Mention software response.}} \\ \hline
\multicolumn{4}{|l|}{\textbf{Comments:}} \\ \hline
\multicolumn{4}{|c|}{\textbf{Passed      Failed     Not Executed}} \\ \hline
\end{tabular}
\label{tab:3}
\end{table}
\section{Test Metrics}
Summarize here the common ground of attributes of test case metrics as explained in Table \ref{tab:4}.
\begin{table}[!ht]
\centering
\caption{Sample Test case Matric.No.1}

\begin{tabular}{|l|l|}
\hline
\textbf{Metric}                      & \textbf{Purpose}                                                                                                        \\ \hline
\textbf{Number of Test Cases}        & \begin{tabular}[c]{@{}l@{}}Total number of test cases that you have\\ developed for  your system.\end{tabular}          \\ \hline
\textbf{Number of Test Cases Passed} & The number of test cases that successfully passed                                                                       \\ \hline
\textbf{Number of Test Cases Failed} & The number of test cases that failed                                                                                    \\ \hline
\textbf{Test Case Defect Density}    & \begin{tabular}[c]{@{}l@{}}(No of test cases failed * 100) \\ No of test cases executed\end{tabular}                    \\ \hline
\textbf{Test Case Effectiveness}     & \begin{tabular}[c]{@{}l@{}}No of defects detected using test cases *100\\ Total number of defects detected\end{tabular} \\ \hline
\textbf{Traceability Matrix} &
  \begin{tabular}[c]{@{}l@{}}Traceability is the ability to determine that each\\  the feature has a source in requirements and each \\ requirement has a corresponding implemented feature.\end{tabular} \\ \hline
\end{tabular}
\label{tab:4}
\end{table}
\subsection{Sample Test case Metric.No.2}
Similar to the above table, add here for test metrics 2,3,4.\\
